\addcontentsline{toc}{section}{Introduction}
\section*{Introduction}

Ce chapitre présente la conception de BioGate. Dans un premier temps, les besoins sont analysés. Ensuite, l'architecture est décrite en distinguant ce qui est visé de ce qui a été réalisé dans le prototype. Le pipeline d'analyse par intelligence artificielle, l'authentification multifacteur, le modèle de données, les alertes à distance, les réglages, la gestion des pannes et les choix techniques sont enfin détaillés.

\section{Analyse des besoins}

Cette section identifie ce que le système doit faire, les contraintes du projet et les acteurs concernés.

\subsection{Besoins fonctionnels}

Le système doit répondre aux besoins suivants :
\begin{itemize}
  \item détecter et suivre les personnes présentes dans le champ de la caméra ;
  \item identifier les personnes enrôlées à partir de leur visage ;
  \item exiger, selon le profil de sécurité choisi, un second facteur : une empreinte digitale ou un test de vivacité ;
  \item adapter les profils enrôlés aux changements d'apparence ;
  \item classer en intrus un visage inconnu, conserver une photo et alerter le propriétaire à distance ;
  \item journaliser les événements et permettre leur consultation ;
  \item permettre l'enrôlement, la configuration et la supervision depuis un tableau de bord.
\end{itemize}

\subsection{Besoins non fonctionnels}

\begin{itemize}
  \item \textbf{Confidentialité} : l'analyse des images se fait en local et les photos d'intrus sont supprimées automatiquement après 72 heures ;
  \item \textbf{Réactivité} : la fenêtre d'identification d'un visage inconnu est d'une seconde par défaut ;
  \item \textbf{Sobriété} : les modèles retenus sont légers, l'image est réduite avant la détection et les vecteurs des profils sont gardés en mémoire pour éviter d'interroger la base à chaque image ;
  \item \textbf{Résilience} : la caméra se reconnecte seule au dernier réseau Wi-Fi enregistré et la purge des photos s'exécute aussi au démarrage du serveur, ce qui rattrape un arrêt prolongé du serveur ;
  \item \textbf{Fiabilité des alertes} : une fenêtre de grâce\glsadd{fenetregrace}, de 60 secondes par défaut, évite d'alerter à tort pour un résident pas encore identifié ;
  \item \textbf{Contrôle des accès au système} : la caméra et le tableau de bord sont protégés par une clé ou par un code de couplage.
\end{itemize}

\subsection{Contraintes}

Le prototype a été réalisé avec des moyens limités. Il s'exécute sur un ordinateur personnel. Aucun lecteur d'empreinte ni actionneur physique (gâche électrique, éclairage) n'a été réalisé : ces éléments sont simulés. Le portage du serveur sur un micro-ordinateur monocarte de type Raspberry Pi est visé à terme, mais il n'a pas été testé et le traitement y serait plus lent.

\subsection{Cahier des charges}\label{sec:cdc}

Le cahier des charges est le document de départ du projet, dont le candidat et l'encadrant sont les signataires. Il fixe le thème, les objectifs, l'architecture visée, les scénarios d'usage et les livrables attendus. Il est distinct de ce mémoire, qui rend compte de ce qui a été réalisé et mesuré. Il est conservé dans le dépôt du projet. Sur sept points, le prototype s'en écarte :
\begin{itemize}
  \item \textbf{Alertes} : le cahier prévoyait WhatsApp. Telegram a été retenu (section~\ref{sec:choix}) ;
  \item \textbf{Serveur} : le cahier prévoyait un Raspberry Pi 4 ou un ordinateur sous Linux. Le prototype tourne sur un ordinateur sous Windows 10 ;
  \item \textbf{Caméra} : le cahier indique la puce RTL8720DN pour la carte BW21-CBV. La fiche technique et le programme de la carte indiquent la puce RTL8735B, retenue dans ce mémoire ;
  \item \textbf{Score de confiance} : le cahier fixait 60~\% après la reconnaissance du visage. Le prototype affiche la confiance réelle de la reconnaissance, et ne fixe que les valeurs de 85~\% et 100~\% après un second facteur (section~\ref{sec:mfa}) ;
  \item \textbf{Base de données} : le cahier prévoyait trois tables, dont des vecteurs en lecture seule. Le prototype en a six, et un vecteur est ajouté à un profil après confirmation par l'empreinte (section~\ref{sec:adapt}) ;
  \item \textbf{Tableau de bord} : le cahier prévoyait un tableau de bord Web servi par le serveur. Le prototype fournit une application de bureau ;
  \item \textbf{Matériel} : la maquette de porte, le lecteur d'empreinte, la gâche, l'éclairage, le réseau fermé et le mini-onduleur du cahier n'ont pas été réalisés. Ils sont simulés ou restent en perspectives.
\end{itemize}

\subsection{Démarche de développement}

Le développement a suivi une démarche itérative et incrémentale, consignée dans un journal de bord de 30 entrées, conservé dans le dépôt du projet. Chaque phase (détection, suivi, reconnaissance, authentification multifacteur, test de vivacité, tableau de bord, alertes, purge des photos) s'est terminée par des essais à la webcam, qui ont révélé des défauts corrigés au fur et à mesure. Ces notes fournissent les raisons des choix présentés dans ce chapitre, et les défauts du tableau~\ref{tab:defauts} y sont consignés.

\subsection{Notation des schémas}

Les schémas suivent la notation \gls{uml}. Dans un diagramme de séquence, une flèche pleine à pointe pleine est un appel synchrone, une flèche pleine à pointe ouverte un message asynchrone et une flèche pointillée une réponse. Un cadre \emph{alt} regroupe des cas exclusifs, chacun précédé de sa condition entre crochets. Dans le diagramme de cas d'utilisation, une flèche pointillée marquée \guillemotleft{} extend \guillemotright{} relie un cas qui en complète un autre. Dans le diagramme d'états, un rond plein marque l'état initial. Dans le schéma d'architecture, un cadre ou une flèche en pointillés désigne un élément simulé.

\subsection{Acteurs et cas d'utilisation}

Cinq acteurs interagissent avec le système :
\begin{itemize}
  \item \textbf{Résident} : personne enrôlée, c'est-à-dire enregistrée par le propriétaire (\glsadd{enrolement}son visage est transformé en vecteurs, rangés dans un profil), qui entre chez elle. Le code et le tableau de bord l'appellent « VIP » ;
  \item \textbf{Inconnu} : personne non enrôlée, présente devant la caméra ;
  \item \textbf{Propriétaire} : personne qui supervise et configure le système depuis le tableau de bord ;
  \item \textbf{Lecteur d'empreinte} : second facteur, simulé dans le prototype ;
  \item \textbf{Service Telegram} : destinataire des alertes.
\end{itemize}
La figure~\ref{fig:usecase} représente les cas d'utilisation. La reconnaissance d'une personne devant la caméra en est le cas central. Le test de vivacité et la validation d'empreinte l'étendent lorsque l'authentification multifacteur est activée. Un visage inconnu déclenche l'alerte d'intrusion, envoyée au service Telegram.

\begin{figure}[H]
\centering
\begin{tikzpicture}[>={Stealth[length=2mm]},
  uc/.style={draw, ellipse, align=center, font=\scriptsize, minimum width=3.7cm, minimum height=1.0cm, inner sep=1pt},
  ext/.style={font=\tiny, fill=white, inner sep=1pt}]
  % acteurs
  \acteur{Res}{0,2.3}{Résident}
  \acteur{Lec}{0,-0.5}{Lecteur d'empreinte\\(simulé)}
  \acteur{Inc}{0,-3.1}{Inconnu}
  \acteur{Pro}{14.2,0.4}{Propriétaire}
  \acteur{Tel}{4.4,-6.3}{Service\\Telegram}
  % système
  \draw (1.9,3.6) rectangle (12.3,-4.2);
  \node[font=\scriptsize\bfseries] at (7.1,3.3) {Système BioGate};
  % cas d'utilisation
  \node[uc] (U1) at (4.4,2.3) {Se faire reconnaître};
  \node[uc] (U2) at (4.4,0.7) {Passer le test\\de vivacité};
  \node[uc] (U3) at (4.4,-0.8) {Valider une empreinte};
  \node[uc] (U4) at (4.4,-3.0) {Déclencher une alerte\\d'intrusion};
  \node[uc] (U5) at (9.8,2.3) {Superviser (flux, journal,\\statistiques)};
  \node[uc] (U6) at (9.8,0.7) {Gérer les profils};
  \node[uc] (U7) at (9.8,-0.9) {Configurer le système\\et les caméras};
  \node[uc] (U8) at (9.8,-2.6) {Purger les photos\\et le journal};
  % associations
  \draw (Res) -- (U1); \draw (Res) -- (U2);
  \draw (Lec) -- (U3);
  \draw (Inc) -- (U4);
  \draw (Pro) -- (U5); \draw (Pro) -- (U6); \draw (Pro) -- (U7); \draw (Pro) -- (U8);
  \draw (Tel) -- (U4);
  % extensions
  \draw[dashed, -{Straight Barb[length=2mm]}] (U2.north) -- node[ext, right] {\guillemotleft{} extend \guillemotright{}} (U1.south);
  \draw[dashed, -{Straight Barb[length=2mm]}] (U3.west) to[bend left=35] node[ext, left] {\guillemotleft{} extend \guillemotright{}} (U1.west);
\end{tikzpicture}
\caption{Diagramme de cas d'utilisation de BioGate}\label{fig:usecase}
\end{figure}

\section{Architecture du système}

Cette section décrit l'organisation de BioGate et l'état de chaque composant dans le prototype.

\subsection{Vue d'ensemble}

BioGate sépare l'acquisition des images, l'intelligence artificielle et la supervision. La caméra envoie son flux au serveur local, qui analyse les images, décide de l'accès, enregistre les événements et alimente le tableau de bord. La figure~\ref{fig:architecture} présente cette organisation, et le tableau~\ref{tab:composants} indique l'état de chaque composant.

\begin{figure}[H]
\centering
\begin{tikzpicture}[>={Stealth[length=2mm]},
  box/.style={draw, rounded corners, align=center, minimum height=1.1cm, inner sep=4pt, font=\footnotesize},
  sim/.style={box, dashed}]
  \node[box, minimum width=3cm] (cam) at (0,0) {Caméra BW21-CBV\\flux MJPEG protégé};
  \node[box, minimum width=4cm] (srv) at (5.6,0) {Serveur local (Edge)\\API FastAPI\\pipeline d'analyse\\score de confiance\\base SQLite};
  \node[box, minimum width=3.6cm] (dash) at (11.4,1.7) {Tableau de bord\\Electron / React};
  \node[box, minimum width=3.6cm] (alert) at (11.4,0) {Alertes\\Telegram};
  \node[sim, minimum width=3.6cm] (iot) at (11.4,-1.7) {Empreinte, gâche, lumière\\(simulés)};
  \draw[->] (cam) -- (srv);
  \draw[<->] (srv) -- (dash);
  \draw[->] (srv) -- (alert);
  \draw[<->, dashed] (srv) -- (iot);
\end{tikzpicture}
\caption{Architecture de BioGate (les éléments en pointillés sont simulés dans le prototype)}\label{fig:architecture}
\end{figure}

\begin{table}[H]
\centering
\caption{Composants du système et état dans le prototype}\label{tab:composants}
\small
\begin{tabularx}{\textwidth}{@{}p{3.1cm}XX@{}}
\toprule
\textbf{Composant} & \textbf{Rôle} & \textbf{État dans le prototype} \\
\midrule
Caméra BW21-CBV & Capture et diffusion du flux MJPEG protégé & Réalisé et testé : flux reçu par le serveur, reconnexion au réseau enregistré après un redémarrage \\
Serveur local & Analyse des images, authentification, base de données, alertes & Réalisé, exécuté sur un ordinateur personnel \\
Tableau de bord & Supervision, enrôlement, configuration, alertes & Réalisé et testé ; le flux de la caméra et les alertes s'y affichent \\
Alertes à distance & Notification d'intrusion avec photo & Réalisé avec Telegram, testé en conditions réelles, avec suppression des messages après 24 heures ; courrier électronique et WhatsApp laissés en perspectives \\
Lecteur d'empreinte & Second facteur d'authentification & Conçu et simulé \\
Gâche et éclairage & Ouverture et éclairage d'accueil ou d'alerte & Conçus et simulés \\
Écran \gls{lcd} sur la caméra & Affichage des consignes du test de vivacité & Prévu, non réalisé \\
\bottomrule
\end{tabularx}
\end{table}

Le serveur est conçu pour s'exécuter à terme sur un micro-ordinateur monocarte de type Raspberry Pi ou sous Linux. Le prototype tourne sur un ordinateur personnel. Trois choix limitent le calcul : l'image est réduite avant la détection, les vecteurs des profils sont gardés en mémoire, et une personne déjà identifiée n'est revérifiée que toutes les 0,8 seconde.

\subsection{Caméra Ai-Thinker BW21-CBV}

La caméra est une carte BW21-CBV de la société Ai-Thinker (figure~\ref{fig:carte}). Selon sa fiche technique, elle est bâtie sur la puce RTL8735B de Realtek et offre un Wi-Fi bi-bande (2,4~GHz ou 5~GHz), le Bluetooth basse consommation 5.1 et une programmation avec l'environnement Arduino~\cite{aithinker_bw21_spec}. Elle a été retenue pour cette connectivité bi-bande. Au laboratoire, la carte s'est connectée en 2,4~GHz, et l'auteur l'a essayée sur les deux bandes lors de ses propres essais, hors journal.

\imageouplace{bw21_carte.jpg}{0.5\textwidth}{Kit Ai-Thinker BW21-CBV. Source de l'image : Youpilab}{fig:carte}

Le micrologiciel diffuse l'image en format \gls{mjpeg}, en 640 par 480 pixels, taille analysée par le serveur. Quatre mesures protègent le flux :
\begin{itemize}
  \item \textbf{Authentification} : l'accès exige un identifiant et une clé de 12 à 32 caractères alphanumériques, choisie par l'auteur. La clé est saisie, avec le réseau Wi-Fi, sur le port série au premier démarrage, puis mémorisée dans la mémoire flash de la carte ;
  \item \textbf{Comparaison en temps constant} : la durée de la réponse ne révèle pas à partir de quel caractère la clé diffère ;
  \item \textbf{Requête bornée} : la lecture est limitée à 1\,024 octets et à 3 secondes, ce qui empêche un client silencieux de bloquer la caméra ;
  \item \textbf{Verrouillage} : après cinq refus de suite, la caméra refuse toute connexion pendant trente secondes.
\end{itemize}
Les échanges ne sont pas chiffrés : la bibliothèque de la carte n'offre ni \gls{https} ni filtrage par adresse. La protection repose donc aussi sur un réseau Wi-Fi chiffré. Les réglages mémorisés dans la carte ne sont pas chiffrés. Au redémarrage, la carte se reconnecte seule au dernier réseau enregistré, ce qui a été vérifié. Si le Wi-Fi tombe, une boucle du programme relance la connexion ; cette reprise n'a pas été éprouvée par une coupure provoquée.

\subsection{Serveur et interface de programmation}\label{sec:routes}

Le serveur est une application Python bâtie avec FastAPI~\cite{fastapi_docs}. Il expose 37 routes de type \gls{api}. Une route\glsadd{route} est une adresse du serveur à laquelle un programme adresse une demande, par exemple la liste des profils. Le jeton\glsadd{jeton} est un code secret que le serveur remet au tableau de bord après le couplage (section~\ref{sec:couplage}) et que celui-ci joint à chaque demande. 26 routes exigent ce jeton. Les 11 autres se répartissent ainsi :
\begin{itemize}
  \item le couplage (deux routes) et la vérification de santé ;
  \item les six routes appelées par le programme client de développement : analyse d'une image, enrôlement, état de l'accès, résultat de l'empreinte, demande et résultat du test de vivacité ;
  \item les deux routes des profils de sécurité (liste et application d'un profil), que le tableau de bord appelle.
\end{itemize}
La protection des deux dernières catégories n'est pas traitée dans le prototype (chapitre~\ref{chap:3}).

Le serveur lit lui-même les caméras (\gls{usb}, flux \gls{mjpeg} ou \gls{rtsp}). Chaque source est ouverte dans un fil d'exécution distinct, et les images sont déposées dans une file de quatre places : si elle est pleine, l'image est abandonnée, ce qui évite l'accumulation de retard. Un seul fil analyse les images, avec un suiveur par caméra. Les cadres et les noms sont dessinés à la cadence de la caméra avec les dernières annotations calculées, si bien que la vidéo reste fluide. Le débit de l'analyse est estimé de deux à cinq images par seconde sur le processeur seul, sans mesure complète. Ce fonctionnement est celui utilisé avec la caméra BW21-CBV et le tableau de bord.

Deux outils servent au développement et ne font pas partie du produit. Le programme client (\texttt{client\_camera.py}) envoie les images de la webcam au serveur, simule l'empreinte, la porte et la lumière par des touches du clavier, et calcule en local les coefficients du test de vivacité, que le serveur accepte sans pouvoir les vérifier. Le programme autonome (\texttt{main.py}) exécute la chaîne d'analyse sans serveur, pour déboguer l'intelligence artificielle seule.

\subsection{Tableau de bord}

Le tableau de bord est une application de bureau développée avec Electron~\cite{electron_docs} et React. Il comporte sept pages :
\begin{itemize}
  \item \textbf{Surveillance} : flux annoté de la caméra, état de chaque individu, score de confiance et progression du test de vivacité ;
  \item \textbf{Journal} : liste des événements, filtrable par période, par type et par nom, exportable au format \gls{csv}, avec un bouton de purge ;
  \item \textbf{VIP} : liste des résidents, enrôlement par fichier image ou par capture de la caméra, retrait d'un profil ;
  \item \textbf{Statistiques} : indicateurs de la période, graphique des accès par jour et liste des résidents les plus actifs ;
  \item \textbf{Configuration} : sources vidéo, profils de sécurité, réglages et état de la purge automatique ;
  \item \textbf{Alertes} : liste des intrusions confirmées, réglage et test des canaux Telegram et courrier électronique, notifications du bureau ;
  \item \textbf{Couplage} : saisie de l'adresse du serveur, puis du code d'accès.
\end{itemize}
Une notification native du système s'affiche lorsqu'une intrusion est confirmée. La figure~\ref{fig:dashboard} montre la page de surveillance.

\imageouplace{dashboard_surveillance.png}{0.95\textwidth}{Page de surveillance du tableau de bord}{fig:dashboard}

\subsection{Sécurité des échanges avec le tableau de bord}\label{sec:couplage}

Le tableau de bord est couplé au serveur comme un appareil domestique :
\begin{enumerate}
  \item le serveur tire au hasard un code à six chiffres et l'affiche dans son terminal ;
  \item l'utilisateur saisit ce code dans le tableau de bord ; il est valable cinq minutes et limité à cinq essais ;
  \item le serveur renvoie un jeton aléatoire, valable trente jours, que l'application conserve chiffré par le système d'exploitation.
\end{enumerate}
Ce jeton est exigé par 26 routes, dont les deux flux vidéo, où il passe dans l'adresse de la requête. Les origines autorisées à appeler le serveur sont limitées à celles du tableau de bord. Les échanges ne sont pas chiffrés : le prototype n'utilise pas HTTPS.

\section{Pipeline de traitement par intelligence artificielle}

Cette section décrit les étapes appliquées à chaque image et la raison de chacune. La figure~\ref{fig:pipeline} les résume.

\begin{figure}[H]
\centering
\begin{tikzpicture}[>={Stealth[length=2mm]},
  box/.style={draw, rounded corners, align=center, minimum height=1.2cm, minimum width=3cm, inner sep=3pt, font=\scriptsize},
  ext/.style={box, dashed}]
  \node[box] (a) at (0,0) {1. Détection\\des personnes\\(YOLOv8n)};
  \node[box] (b) at (3.6,0) {2. Suivi\\par centroïde\\et couleurs};
  \node[box] (c) at (7.2,0) {3. Filtre de visage\\(Haar, puis\\second avis MediaPipe)};
  \node[box] (d) at (10.8,0) {4. Contraste\\(CLAHE)};
  \node[box] (e) at (10.8,-2.2) {5. Vecteur de 512\\nombres\\(FaceNet512)};
  \node[box] (f) at (7.2,-2.2) {6. Comparaison\\(distance cosinus,\\seuil 0,35)};
  \node[box] (g) at (3.6,-2.2) {7. Décision\\et revérification};
  \node[ext] (h) at (0,-2.2) {Score, journal,\\photo, alerte};
  \draw[->] (a) -- (b); \draw[->] (b) -- (c); \draw[->] (c) -- (d);
  \draw[->] (d) -- (e); \draw[->] (e) -- (f); \draw[->] (f) -- (g); \draw[->] (g) -- (h);
\end{tikzpicture}
\caption{Chaîne de traitement d'une image (pas de visage de face : l'analyse s'arrête à l'étape~3)}\label{fig:pipeline}
\end{figure}

\begin{enumerate}
  \item \textbf{Détection des personnes} : le modèle \gls{yolo}v8 nano repère les silhouettes humaines~\cite{ultralytics_yolov8docs}. Il est léger : environ 15 images par seconde sur le processeur d'un ordinateur sous Windows (1\,800 images en 118 secondes). L'image est réduite à 320 pixels pour alléger le calcul, et le seuil de confiance est de 0,5~\cite{ultralytics_predict} : une détection dont le modèle n'est pas sûr à 50~\% est écartée. Deux boîtes sur une même personne, corps et tête, créaient des intrus fantômes. Le recouvrement de deux boîtes se mesure par le rapport entre leur surface commune et leur surface réunie : au-dessus de 0,35, le modèle les fusionne, et une seconde fusion à 0,30 traite les boîtes adjacentes.
  \item \textbf{Suivi} : un suivi par centroïde donne un identifiant stable à chaque individu. Il garde en mémoire une personne disparue pendant 60 images, soit environ deux secondes, contre 20 au départ. Il refuse de rattacher une détection située à plus de 150 pixels, sauf si les couleurs des vêtements concordent. La signature de couleur est un histogramme de teinte (espace \gls{hsv}, 16 classes), et la corrélation, qui mesure la ressemblance de deux histogrammes entre $-1$ et 1, doit dépasser 0,85. Une personne cachée un instant derrière un obstacle garde ainsi son identité, sans nouvelle analyse du visage.
  \item \textbf{Filtre de visage de face} : le modèle lourd n'est pas appelé sur une personne de dos, de profil ou la tête baissée. Une cascade de Haar~\cite{viola2001rapid,opencv_cascade} cherche le visage sur une image en niveaux de gris à histogramme égalisé (échelle de 1,1, cinq voisins, taille minimale de 30 par 30 pixels). Elle prenait le col et un bouton d'une chemise pour un visage : un second avis du modèle de points du visage de MediaPipe écarte ces faux visages. Sur les essais, la chemise est rejetée et les trois portraits de test sont acceptés.
  \item \textbf{Correction du contraste} : la méthode \gls{clahe}~\cite{zuiderveld1994clahe,opencv_clahe} est appliquée à la luminosité du visage, un des trois canaux de l'espace de couleur LAB. Elle rétablit les traits en contre-jour et stabilise la distance mesurée d'une image à l'autre.
  \item \textbf{Représentation du visage} : le modèle FaceNet512, de la bibliothèque DeepFace, transforme le visage aligné (yeux redressés) en un vecteur de 512 nombres~\cite{schroff2015facenet,serengil2020lightface}. Le FaceNet décrit au chapitre~\ref{chap:1} produit 128 nombres ; DeepFace en propose une variante à 512, que BioGate emploie. Il a été choisi après l'essai de VGG-Face avec une détection basique, trop imprécis sous la lumière d'une webcam.
  \item \textbf{Comparaison} : le vecteur est comparé à ceux des profils par la distance cosinus\glsadd{distancecosinus}. Cette distance vaut 0 pour deux vecteurs identiques et croît avec leur différence. Une personne est reconnue si la distance au profil le plus proche est inférieure au seuil de 0,35. La confiance affichée vaut $(1 - d) \times 100$, où $d$ est la distance : elle est de 65~\% au seuil. Une distance de 0,22 donne par exemple une confiance de 78~\% et une reconnaissance, une distance de 0,40 une confiance de 60~\% et un refus. Les vecteurs sont gardés en mémoire, car interroger la base à chaque image ralentissait le système. Ils sont chargés au démarrage et rechargés après un enrôlement ou un apprentissage.
  \item \textbf{Décision et journalisation} :
    \begin{itemize}
      \item un résident est accepté dès la première correspondance sous le seuil ;
      \item un visage inconnu est jugé une seconde après son apparition ;
      \item une personne identifiée est revérifiée toutes les 0,8 seconde, pour qu'un résident ne soit pas remplacé par un intrus ou masqué après sa validation ;
      \item trois échecs de suite sont nécessaires pour la remettre en question, afin qu'un changement d'expression ou de lumière ne crée pas un faux intrus ;
      \item un résident qui retombe en analyse dispose de cinq secondes au lieu d'une ;
      \item un intrus provisoire reconnu comme résident à une distance inférieure à 0,25 est reclassé aussitôt, et son alerte annulée.
    \end{itemize}
    L'événement est enregistré dans la base.
  \item \textbf{Photo de l'intrus présumé} : lorsqu'un individu devient intrus provisoire, une photo de sa silhouette est rattachée à l'événement et à l'alerte. Elle entre dans la purge de 72 heures. Seul un individu soupçonné est photographié, jamais un résident reconnu. Le programme client de développement enregistre à l'apparition la silhouette de chaque individu, résidents compris ; la purge supprime celles qui ne sont rattachées à aucun événement.
\end{enumerate}

La bibliothèque DeepFace fixe le seuil par défaut de FaceNet512 à 0,30~\cite{serengil_deepface_thresholds}. BioGate retient 0,35, valeur calibrée de façon empirique : 0,20 refusait trop de résidents et 0,60 acceptait trop d'inconnus. Relever le seuil réduit les refus de résidents mais augmente le risque de confondre deux personnes ; l'abaisser produit l'effet inverse. C'est le compromis entre les taux de fausse acceptation et de faux rejet défini au chapitre~\ref{chap:1}.

Un intrus est un visage inconnu, et non une silhouette sans visage. Une silhouette sans visage reste en attente de visage. Si elle disparaît après plus de trois secondes de présence sans avoir montré de visage, elle laisse seulement une trace dans le journal (passage sans visage), sans photo ni alerte. Un autre réglage, la durée au bout de laquelle une silhouette sans visage devient intrus provisoire, est à « jamais » dans les trois premiers profils du tableau~\ref{tab:profils} et vaut 5 secondes en haute sécurité. Ce choix répond à des faux intrus observés sur une main et sur un torse.

\section{Authentification multifacteur et score de confiance}\label{sec:mfa}

Cette section décrit les états d'un individu, le score de confiance, le second facteur, le test de vivacité et l'adaptation des profils.

\subsection{Machine à états}

Chaque individu suivi est géré par une machine à états (figure~\ref{fig:etats}). Ses états sont les suivants :
\begin{itemize}
  \item \textbf{Repos} : état initial, avant l'arrivée d'un individu ;
  \item \textbf{Analyse} : l'individu vient d'apparaître et son visage est en cours d'analyse ;
  \item \textbf{Accès accordé} : un résident est reconnu et, si besoin, son second facteur est validé ;
  \item \textbf{Second facteur en attente} : le visage est reconnu et l'authentification multifacteur est activée ;
  \item \textbf{Accès refusé} : le second facteur a échoué ou le délai est dépassé ;
  \item \textbf{Intrus provisoire} : le visage est resté inconnu, la fenêtre de grâce est ouverte ;
  \item \textbf{Intrus confirmé} : la fenêtre de grâce s'est terminée sans que l'individu soit reconnu comme résident.
\end{itemize}

\begin{figure}[H]
\centering
\begin{tikzpicture}[>={Straight Barb[length=2mm]},
  st/.style={draw, rounded corners, align=center, minimum height=1cm, inner sep=3pt, font=\scriptsize},
  lab/.style={font=\tiny, align=center, fill=white, inner sep=1pt}]
  \node[circle, fill, inner sep=2.5pt] (init) at (0,0) {};
  \node[st, minimum width=1.6cm] (idle) at (2.0,0) {IDLE};
  \node[st, minimum width=2.4cm] (ana) at (4.9,0) {ANALYZING};
  \node[st, minimum width=3.6cm] (grant) at (10.0,2.6) {GRANTED\\accès accordé};
  \node[st, minimum width=3.6cm] (fp) at (10.0,0) {FINGERPRINT\_PENDING\\second facteur};
  \node[st, minimum width=2.6cm] (den) at (14.9,0) {DENIED\\accès refusé};
  \node[st, minimum width=3.6cm] (ip) at (10.0,-2.6) {INTRUDER\_PENDING\\intrus provisoire};
  \node[st, minimum width=3.6cm] (ic) at (14.9,-2.6) {INTRUDER\_CONFIRMED\\alerte envoyée};
  \draw[->] (init) -- (idle);
  \draw[->] (idle) -- node[lab, above] {individu\\détecté} (ana);
  \draw[->] (ana.north) -- node[lab, above left] {visage reconnu\\{[}MFA désactivée{]}} (grant.west);
  \draw[->] (ana) -- node[lab, above] {visage reconnu\\{[}MFA activée{]}} (fp);
  \draw[->] (fp) -- node[lab, right] {empreinte ou\\vivacité réussie} (grant);
  \draw[->] (fp) -- node[lab, above] {échec ou\\{[}15 s{]}} (den);
  \draw[->] (ana.south) -- node[lab, above right, pos=0.55] {visage inconnu\\{[}après 1 s{]}} (ip.west);
  \draw[->] (ip) -- node[lab, above] {grâce\\écoulée} (ic);
  \draw[->] (ip.south) -- ++(0,-0.7) -| node[lab, pos=0.25, below] {résident reconnu / alerte annulée} (ana.south west);
\end{tikzpicture}
\caption{Diagramme d'états d'un individu suivi}\label{fig:etats}
\end{figure}

Le score de confiance accompagne chaque état :
\begin{itemize}
  \item \textbf{Visage seul} : confiance de la reconnaissance, $(1 - d) \times 100$, recalculée à chaque revérification ;
  \item \textbf{Second facteur validé} : valeur fixe, 85~\% après un test de vivacité réussi et 100~\% après une empreinte validée.
\end{itemize}
Ce pourcentage mesure la proximité de deux vecteurs et non la probabilité que la personne soit la bonne.

\subsection{Activer et utiliser le second facteur}

Un seul réglage commande le second facteur : \emph{MFA obligatoire}, dans la page de configuration. Il active l'\gls{mfa}. Il est désactivé dans le profil par défaut (domicile) et activé dans les profils bureau et haute sécurité. Le tableau~\ref{tab:mfa} compare les deux états.

\begin{table}[H]
\centering
\caption{Fonctionnement selon que le second facteur est exigé ou non}\label{tab:mfa}
\small
\begin{tabularx}{\textwidth}{@{}>{\raggedright\arraybackslash}p{3.3cm}>{\raggedright\arraybackslash}X>{\raggedright\arraybackslash}X@{}}
\toprule
 & \textbf{MFA désactivée (profil par défaut)} & \textbf{MFA activée} \\
\midrule
Décision d'accès & Dès que le visage est reconnu & Après la reconnaissance du visage, puis la réussite d'un second facteur \\
Action de la personne & Aucune & Exécuter le geste affiché au tableau de bord, ou poser le doigt sur le lecteur (simulé dans le prototype) ; le premier qui réussit suffit \\
Délais & Aucun & 15 secondes pour l'empreinte, 10 pour le geste \\
Score affiché & Confiance du visage, calculée & 85~\% après le geste, 100~\% après l'empreinte \\
Enrichissement du profil & Impossible & Possible : le vecteur du visage est ajouté si l'empreinte confirme \\
\bottomrule
\end{tabularx}
\end{table}

Lorsque le second facteur est exigé, l'empreinte et le test de vivacité sont deux voies concurrentes : la première qui réussit accorde l'accès.
\begin{itemize}
  \item \textbf{Empreinte} : elle dispose de 15 secondes. Passé ce délai, ou si l'empreinte est refusée, l'accès est refusé ;
  \item \textbf{Vivacité} : elle dispose de 10 secondes. Si elle échoue ou expire, ce canal se ferme sans refuser l'accès, et l'empreinte reste possible.
\end{itemize}
Le réglage \emph{vivacité activée} retire le test de vivacité du second facteur : l'empreinte est alors la seule voie, comme dans le profil haute sécurité. Il prend effet sans redémarrage du serveur.

\subsection{Test de vivacité}

Le test de vivacité vérifie qu'une personne est présente devant la caméra. Il s'appuie sur le modèle de points du visage de MediaPipe, qui fournit 52 coefficients d'expression compris entre 0 et 1~\cite{google_facelandmarker,mediapipe_blendshape_card}. Deux types de coefficients sont utilisés : la fermeture des yeux et le sourire, moyennés sur les deux côtés. À chaque tentative, le système tire au sort un seul geste :
\begin{itemize}
  \item \textbf{Clignements} : deux à quatre clignements, un clignement étant compté lorsque la fermeture des yeux dépasse 0,35 ;
  \item \textbf{Sourire} : un sourire tenu une à deux secondes, dont le coefficient doit dépasser 0,40. La mesure ne repart de zéro que si le coefficient redescend sous 0,25.
\end{itemize}
La personne dispose de dix secondes. Le seuil du sourire a été abaissé de 0,45 à 0,40 après des séances réelles, où les pics atteignent 0,60 à 0,75.

Le tirage au sort rend le geste imprévisible : une photographie ne peut exécuter aucun geste, et une vidéo préparée doit contenir exactement le geste demandé. Ce principe de défi tiré au hasard est celui de Widjaya et Wicaksana, qui l'appliquent à d'autres gestes (ouvrir la bouche, cligner des yeux, tourner la tête)~\cite{widjaya2023challenge}. Une première version exigeait une validation faciale ininterrompue de cinq à six secondes. Elle a été écartée, car une photo immobile la franchirait plus facilement.

BioGate n'utilise qu'une caméra ordinaire en deux dimensions, sans capteur de profondeur ni infrarouge. Il ne peut pas contrôler le relief du visage, ce que fait par exemple la caméra TrueDepth de Face ID (chapitre~\ref{chap:1}). La protection contre une photographie repose donc sur le seul défi actif, dont l'efficacité n'a pas été mesurée (chapitre~\ref{chap:3}). L'ajout d'un capteur de profondeur est une perspective.

\subsection{Adaptation des profils par apprentissage continu}\label{sec:adapt}

L'enrôlement initial d'une personne est une action volontaire du propriétaire. Ensuite, le profil s'enrichit en trois étapes :
\begin{enumerate}
  \item pendant l'attente du second facteur, le système met de côté le vecteur du visage courant ;
  \item si l'empreinte confirme l'identité, ce vecteur est ajouté au profil ;
  \item chaque profil conserve au plus 30 vecteurs, les plus récents.
\end{enumerate}
Cet ajout n'a lieu que si le second facteur est exigé et qu'il confirme l'identité. Une première version ajoutait des vecteurs sans supervision. Elle a été abandonnée, car le modèle pouvait apprendre de ses propres erreurs et dériver. Les vecteurs des intrus ne sont jamais conservés.

Le lecteur d'empreinte n'existe pas dans le prototype : la confirmation est simulée par une commande envoyée à la route correspondante du serveur. Le mécanisme est donc éprouvé par simulation, non avec un capteur réel.

\subsection{Scénarios d'interaction}

Trois diagrammes de séquence détaillent les échanges, à partir du code du serveur.

Le premier (figure~\ref{fig:seq1}) concerne un résident qui se présente devant la caméra. Le serveur détecte la personne, calcule le vecteur de son visage et le compare à ceux des profils. Si la distance est inférieure au seuil, l'événement est enregistré et le tableau de bord reçoit l'identité et le score. Lorsque l'authentification multifacteur est activée, le tableau de bord affiche le geste demandé, que la personne exécute devant la caméra.

\begin{figure}[H]
\centering
\begin{tikzpicture}[>={Stealth[length=2mm]}]
  \seqhead{A}{0}{Résident}
  \seqhead{B}{2.9}{Caméra}
  \seqhead{C}{6.0}{Serveur}
  \seqhead{D}{9.0}{Base\\SQLite}
  \seqhead{E}{12.0}{Tableau\\de bord}
  \seqlife{A}{-9.9}\seqlife{B}{-9.9}\seqlife{C}{-9.9}\seqlife{D}{-9.9}\seqlife{E}{-9.9}
  \seqactiv{6.0}{-2.0}{-5.4}
  \seqactiv{9.0}{-4.0}{-4.5}
  \seqactiv{6.0}{-6.5}{-9.2}
  \seqasync{0}{2.9}{-1.3}{se présente de face}
  \seqasync{2.9}{6.0}{-2.0}{images}
  \seqself{6.0}{-2.6}{identification du visage (distance $<$ 0,35)}{2.2cm}
  \seqsync{6.0}{9.0}{-4.0}{journalise VIP\_ENTRY}
  \seqret{9.0}{6.0}{-4.5}{ok}
  \seqasync{6.0}{12.0}{-5.1}{identité et score de confiance}
  \draw[dashed] (-0.7,-5.8) rectangle (12.9,-9.6);
  \node[draw, fill=white, font=\tiny, anchor=north west] at (-0.7,-5.8) {alt {[}MFA activée{]}};
  \seqasync{6.0}{12.0}{-6.5}{geste demandé}
  \seqasync{0}{2.9}{-7.1}{exécute le geste}
  \seqasync{2.9}{6.0}{-7.7}{images}
  \seqself{6.0}{-8.1}{geste vérifié (85 \%)}{2.2cm}
  \seqasync{6.0}{12.0}{-9.0}{accès accordé}
\end{tikzpicture}
\caption{Diagramme de séquence : reconnaissance d'un résident}\label{fig:seq1}
\end{figure}

Le deuxième (figure~\ref{fig:seq2}) concerne un visage inconnu. Une seconde après l'apparition du visage, l'individu devient un intrus provisoire et une photo est prise. Pendant la fenêtre de grâce, la reconnaissance d'un résident annule l'alerte et supprime la photo. Sinon, l'intrus est confirmé et le message part vers Telegram avec la photo.

\begin{figure}[H]
\centering
\begin{tikzpicture}[>={Stealth[length=2mm]}]
  \seqhead{A}{0}{Inconnu}
  \seqhead{B}{2.7}{Caméra}
  \seqhead{C}{5.5}{Serveur}
  \seqhead{D}{8.3}{Base\\SQLite}
  \seqhead{E}{10.9}{Service\\Telegram}
  \seqhead{F}{13.5}{Propriétaire}
  \seqlife{A}{-11.8}\seqlife{B}{-11.8}\seqlife{C}{-11.8}\seqlife{D}{-11.8}\seqlife{E}{-11.8}\seqlife{F}{-11.8}
  \seqactiv{5.5}{-2.0}{-4.8}
  \seqactiv{8.3}{-4.2}{-4.7}
  \seqactiv{5.5}{-6.0}{-6.9}
  \seqactiv{5.5}{-7.6}{-8.8}
  \seqactiv{5.5}{-9.9}{-11.2}
  \seqactiv{8.3}{-9.9}{-10.3}
  \seqasync{0}{2.7}{-1.3}{se présente}
  \seqasync{2.7}{5.5}{-2.0}{images}
  \seqself{5.5}{-2.6}{visage inconnu après 1 s : intrus provisoire, photo}{2.0cm}
  \seqsync{5.5}{8.3}{-4.2}{journalise INTRUDER\_PENDING}
  \seqret{8.3}{5.5}{-4.7}{ok}
  \draw[dashed] (-0.7,-5.2) rectangle (14.4,-11.5);
  \node[draw, fill=white, font=\tiny, anchor=north west] at (-0.7,-5.2) {alt {[}résident reconnu pendant la grâce{]}};
  \seqself{5.5}{-6.0}{alerte annulée, photo supprimée}{2.0cm}
  \draw[dashed] (-0.7,-7.2) -- (14.4,-7.2);
  \node[font=\tiny, anchor=north west] at (-0.6,-7.2) {{[}individu hors du champ{]}};
  \seqself{5.5}{-7.6}{minuteur annulé, photo supprimée, aucune alerte}{2.0cm}
  \draw[dashed] (-0.7,-9.3) -- (14.4,-9.3);
  \node[font=\tiny, anchor=north west] at (-0.6,-9.3) {{[}sinon, grâce écoulée{]}};
  \seqsync{5.5}{8.3}{-9.9}{INTRUDER\_CONFIRMED}
  \seqasync{5.5}{10.9}{-10.5}{message et photo}
  \seqasync{10.9}{13.5}{-11.1}{notification}
\end{tikzpicture}
\caption{Diagramme de séquence : détection d'un visage inconnu et alerte}\label{fig:seq2}
\end{figure}

Le troisième (figure~\ref{fig:seq3}) concerne l'enrichissement d'un profil. Si l'empreinte est validée, le vecteur mis de côté est ajouté au profil, puis la mémoire des vecteurs est rechargée.

\begin{figure}[H]
\centering
\begin{tikzpicture}[>={Stealth[length=2mm]}]
  \seqhead{A}{0}{Résident}
  \seqhead{B}{3.3}{Serveur}
  \seqhead{C}{6.8}{Lecteur d'empreinte\\(simulé)}
  \seqhead{D}{10.2}{Base\\SQLite}
  \seqlife{A}{-7.4}\seqlife{B}{-7.4}\seqlife{C}{-7.4}\seqlife{D}{-7.4}
  \seqactiv{3.3}{-1.9}{-6.6}
  \seqactiv{6.8}{-3.2}{-4.4}
  \seqactiv{10.2}{-5.0}{-5.5}
  \seqasync{0}{3.3}{-1.3}{visage présenté}
  \seqself{3.3}{-1.9}{visage reconnu : vecteur mis de côté}{2.3cm}
  \seqsync{3.3}{6.8}{-3.2}{attend le doigt}
  \seqasync{0}{6.8}{-3.8}{pose le doigt}
  \seqret{6.8}{3.3}{-4.4}{empreinte validée}
  \seqsync{3.3}{10.2}{-5.0}{ajoute le vecteur (30 au plus)}
  \seqret{10.2}{3.3}{-5.5}{ok}
  \seqself{3.3}{-6.0}{recharge les vecteurs en mémoire}{2.3cm}
\end{tikzpicture}
\caption{Diagramme de séquence : apprentissage continu confirmé par l'empreinte}\label{fig:seq3}
\end{figure}

\section{Modèle de données et protection de la vie privée}

Cette section décrit la base de données et les mesures qui limitent la conservation des données.

La base est un fichier SQLite~\cite{sqlite_about} local, exclu du dépôt de code. Elle contient six tables (figure~\ref{fig:donnees}). Un profil peut avoir plusieurs vecteurs, au plus 30, et plusieurs événements au journal. Les tables de configuration, des caméras et des messages Telegram sont indépendantes des autres.

\begin{figure}[H]
\centering
\begin{tikzpicture}[tbl/.style={draw, rectangle split, rectangle split parts=2, font=\scriptsize, align=left, minimum width=4.2cm, rectangle split part fill={gray!15, white}}]
  \node[tbl] (P) at (0,0) {\textbf{profiles}\nodepart{two}\underline{id}\\name (unique)\\role (VIP ou INTRUS)\\created\_at};
  \node[tbl] (E) at (-5.6,0) {\textbf{embeddings}\nodepart{two}\underline{id}\\profile\_id (FK)\\vector\_json (512 nombres)\\created\_at};
  \node[tbl] (L) at (5.6,0) {\textbf{event\_logs}\nodepart{two}\underline{id}\\profile\_id (FK)\\confidence\\event\_type\\snapshot\_path\\timestamp};
  \node[tbl] (C) at (-5.6,-4) {\textbf{cameras}\nodepart{two}\underline{id}\\cam\_id (unique)\\name, type (usb, mjpeg, rtsp)\\url, usb\_index, zone\\active, created\_at};
  \node[tbl] (S) at (0,-4) {\textbf{system\_config}\nodepart{two}\underline{key}\\value\\updated\_at};
  \node[tbl] (T) at (5.6,-4) {\textbf{telegram\_messages}\nodepart{two}\underline{id}\\chat\_id\\message\_id\\sent\_at};
  \draw (E.east) -- node[above, pos=0.12, font=\tiny] {0..*} node[above, pos=0.88, font=\tiny] {1} (P.west);
  \draw (L.west) -- node[above, pos=0.12, font=\tiny] {0..*} node[above, pos=0.88, font=\tiny] {1} (P.east);
\end{tikzpicture}
\caption{Modèle de données de BioGate (tables SQLite ; \underline{souligné} : clé primaire, FK : clé étrangère)}\label{fig:donnees}
\end{figure}

Le journal enregistre dix types d'événements : entrée d'un résident, accès refusé, intrus provisoire, intrus confirmé, test de vivacité réussi, empreinte acceptée, empreinte refusée, alerte de fuite, passage sans visage et simple détection. Dans le code, ces types portent des noms en anglais, comme VIP\_ENTRY pour l'entrée d'un résident ou INTRUDER\_CONFIRMED pour un intrus confirmé. Deux types demandent une précision :
\begin{itemize}
  \item \textbf{Alerte de fuite} : un visage a été vu, puis l'individu a disparu avant le verdict ;
  \item \textbf{Passage sans visage} : une silhouette est restée sans jamais montrer de visage.
\end{itemize}
Un intrus reçoit une ligne dans la table des profils, sous un nom comme « Intrus\_001 », pour que le journal puisse s'y rattacher. Son vecteur n'est jamais conservé. Les vecteurs des résidents ne sont pas chiffrés dans la base (chapitre~\ref{chap:3}).

Plusieurs mesures limitent la conservation des données :
\begin{itemize}
  \item les photos sont supprimées automatiquement après 72 heures, au démarrage du serveur puis toutes les heures ; seul le chemin de la photo est effacé, les lignes du journal restent pour l'historique et les statistiques ;
  \item la photo d'un intrus provisoire qui ne sera jamais confirmé est supprimée immédiatement ;
  \item les messages d'alerte envoyés par Telegram sont supprimés de la conversation après 24 heures par défaut ;
  \item un événement du journal peut être supprimé avec sa photo ;
  \item les identifiants des canaux d'alerte ne sont jamais renvoyés en clair par l'interface, mais ils sont enregistrés en clair dans la base, protégée seulement par le système de fichiers.
\end{itemize}
Cette durée limitée de conservation des photos répond au principe de limitation de la conservation du \gls{rgpd}~\cite{ue2016rgpd} et aux articles 383 et 433 du Code du numérique béninois~\cite{apdp2018codenumerique}, vus au chapitre~\ref{chap:1}.

\section{Alertes à distance}

Cette section décrit la confirmation d'un intrus et l'envoi de l'alerte. Le déroulement est le suivant :
\begin{enumerate}
  \item un visage inconnu est classé intrus provisoire, et une fenêtre de grâce s'ouvre ;
  \item si un résident est reconnu pendant la fenêtre, l'alerte est annulée en silence et la photo supprimée ;
  \item si l'individu quitte le champ de la caméra pendant la fenêtre, le minuteur est annulé, la photo supprimée et aucune alerte n'est envoyée ; seul l'événement d'intrus provisoire reste au journal ;
  \item sinon, l'intrus est confirmé et le serveur envoie un message avec la photo par Telegram.
\end{enumerate}
La fenêtre de grâce dure 60 secondes dans le profil par défaut, et de 5 à 120 secondes selon le profil. Un délai minimal de 30 secondes sépare deux alertes, pour éviter plusieurs messages lorsque le suivi attribue plusieurs identités à la même personne. À l'installation, l'envoi est désactivé (réglage du fichier de configuration). Appliquer le profil portail, domicile ou haute sécurité active Telegram et le courrier électronique (tableau~\ref{tab:profils}), mais rien n'est envoyé tant que le jeton du robot et l'identifiant de la conversation, ou les réglages du courrier, n'ont pas été saisis dans le tableau de bord. Le troisième cas est une limite : un intrus qui quitte le champ avant la fin de la fenêtre n'est pas signalé (chapitre~\ref{chap:3}).

Le courrier électronique et WhatsApp sont des perspectives. Le module d'envoi par courrier est fonctionnel : il joint la photo au message, après négociation d'une connexion chiffrée~\cite{rfc3207}, et ses tests automatiques utilisent un serveur simulé. Il n'a pas été essayé avec un vrai serveur de messagerie.

La photo envoyée par Telegram quitte le traitement local et transite par un service extérieur. Les messages envoyés sont donc supprimés de la conversation après un délai, de 24 heures par défaut :
\begin{itemize}
  \item chaque message envoyé est enregistré dans une table, avec l'identifiant de la conversation et celui du message ;
  \item la tâche de purge, au démarrage du serveur puis toutes les heures, supprime ceux qui ont dépassé le délai ;
  \item le délai est plafonné à 47 heures, car l'interface des robots ne supprime qu'un message de moins de 48 heures~\cite{telegram_botapi} ;
  \item une route protégée par le jeton permet de lancer cette suppression à la main.
\end{itemize}
Un essai réel a supprimé un message de test. Le chemin d'une alerte avec photo n'est couvert que par des tests automatiques, et les messages envoyés avant cette version restent dans la conversation.

\section{Paramètres et profils de sécurité}

Cette section présente les réglages du système et les profils de sécurité. Les réglages sont enregistrés dans la base et rechargés au démarrage du serveur. Quatre profils regroupent des valeurs prêtes à l'emploi : portail, domicile (profil par défaut), bureau et haute sécurité. Chaque réglage reste modifiable. Au premier démarrage, les valeurs du fichier de configuration s'appliquent : elles sont celles du profil domicile, sauf les alertes à distance, désactivées. Les valeurs de 72 heures pour la conservation des photos viennent du cahier des charges (section~\ref{sec:cdc}) ; les autres ont été réglées pendant les essais, sans étude de sensibilité. Le tableau~\ref{tab:profils} donne les principaux réglages des profils.

\begin{table}[H]
\centering
\caption{Principaux réglages des profils de sécurité}\label{tab:profils}
\small
\begin{tabularx}{\textwidth}{@{}>{\raggedright\arraybackslash}Xcccc@{}}
\toprule
\textbf{Réglage} & \textbf{Portail} & \textbf{Domicile} & \textbf{Bureau} & \textbf{Haute sécurité} \\
\midrule
Fenêtre de grâce avant l'alerte (s) & 120 & 60 & 30 & 5 \\
Authentification multifacteur & non & non & oui & oui \\
Test de vivacité dans le second facteur & oui & oui & oui & non \\
Alerte Telegram activée & oui & oui & non & oui \\
Alerte par courrier activée & oui & oui & non & oui \\
Durée d'analyse d'un visage inconnu (s) & 1,5 & 1 & 1 & 0,8 \\
Échecs avant remise en question & 5 & 3 & 3 & 1 \\
Délai de l'empreinte (s) & 20 & 15 & 15 & 10 \\
Conservation des photos (h) & 72 & 72 & 24 & 12 \\
Silhouette sans visage classée intrus après (s) & jamais & jamais & jamais & 5 \\
\bottomrule
\end{tabularx}
\end{table}

Le tableau~\ref{tab:reglages} décrit les réglages principaux : leur rôle, leur valeur par défaut et la façon de les modifier (un curseur ou un interrupteur du tableau de bord, le profil de sécurité, ou le fichier de configuration). La durée d'analyse et l'intervalle de revérification ne prennent effet qu'au redémarrage du serveur.

\begin{table}[H]
\centering
\caption{Principaux réglages}\label{tab:reglages}
\small
\begin{tabularx}{\textwidth}{@{}>{\raggedright\arraybackslash}p{3.2cm}>{\raggedright\arraybackslash}X>{\raggedright\arraybackslash}p{1.3cm}>{\raggedright\arraybackslash}p{2.3cm}@{}}
\toprule
\textbf{Réglage} & \textbf{Rôle, et effet d'une hausse} & \textbf{Défaut} & \textbf{Où} \\
\midrule
Seuil de reconnaissance & Distance maximale entre deux vecteurs pour les juger identiques. Une hausse reconnaît plus de résidents mais augmente le risque de confondre deux personnes. & 0,35 & Curseur \\
Intervalle de revérification & Délai entre deux contrôles d'une personne identifiée. Plus long : moins de calcul, mais un remplacement de personne est détecté plus tard. & 0,8 s & Curseur \\
Durée d'analyse d'un visage inconnu & Temps laissé à la reconnaissance avant de classer l'individu intrus provisoire. Plus long : moins de faux intrus, alerte plus tardive. & 1 s & Curseur \\
MFA obligatoire & Exige un second facteur après la reconnaissance du visage. & non & Interrupteur, profil \\
Vivacité activée & Inclut le test de vivacité dans le second facteur. Désactivé, seule l'empreinte reste possible. & oui & Interrupteur, profil \\
Délai de l'empreinte & Temps accordé pour valider l'empreinte, au-delà duquel l'accès est refusé. & 15 s & Curseur \\
Fenêtre de grâce & Délai avant de confirmer un intrus provisoire. Plus long : plus de temps pour reconnaître un résident, mais alerte plus tardive. & 60 s & Curseur \\
Conservation des photos & Durée avant la suppression automatique des photos. & 72 h & Curseur \\
Suppression des messages Telegram & Délai avant la suppression d'un message d'alerte dans la conversation. Plafonné à 47 h ; 0 désactive. & 24 h & Fichier \\
\bottomrule
\end{tabularx}
\end{table}

\section{Gestion des erreurs et des pannes}

Cette section décrit le comportement du système face aux défaillances. L'échec d'un élément périphérique ne doit pas bloquer l'analyse. Le comportement prévu face aux principales défaillances, tel qu'il est écrit dans le code, est le suivant :

\begin{itemize}
  \item \textbf{Flux vidéo interrompu} : le serveur libère la source et tente de la rouvrir, toutes les secondes après une lecture ratée et toutes les trois secondes si l'ouverture échoue. La caméra BW21-CBV se reconnecte de son côté au dernier réseau enregistré ;
  \item \textbf{Analyse plus lente que la caméra} : les images attendent dans une file de quatre places. Si elle est pleine, l'image est abandonnée, et le fil de capture ne se bloque jamais ;
  \item \textbf{Second avis de MediaPipe indisponible} : un avertissement est journalisé et le détecteur de Haar décide seul. Le risque de prendre un pli de vêtement pour un visage réapparaît alors ;
  \item \textbf{Envoi d'alerte impossible} : l'envoi se fait dans un fil séparé, qui ne ralentit pas l'analyse. L'erreur est journalisée avec un message explicite, dans lequel le jeton Telegram est masqué. Il n'y a pas de nouvelle tentative, et si la photo manque, le message part sans elle ;
  \item \textbf{Second facteur non fourni à temps} : l'accès est refusé après 15 secondes. L'échec du seul test de vivacité ferme ce canal sans refuser l'accès ;
  \item \textbf{Redémarrage du serveur} : les caméras et la configuration enregistrées sont rechargées, et la purge des photos s'exécute au démarrage pour rattraper l'arrêt.
\end{itemize}

Trois pannes ne sont pas traitées :
\begin{itemize}
  \item \textbf{Coupure d'électricité} : le serveur s'arrête. Les états de suivi, les tests de vivacité en cours et le jeton du tableau de bord, tenus en mémoire, sont perdus, et un nouveau couplage est nécessaire ;
  \item \textbf{Échec d'envoi d'une alerte} : le message n'est pas renvoyé ;
  \item \textbf{Perte d'Internet} : seule l'alerte à distance est touchée, et elle n'est pas mise en file pour un envoi ultérieur.
\end{itemize}
Ces limites sont reprises au chapitre~\ref{chap:3}.

\section{Justification des choix techniques}\label{sec:choix}

Cette section présente les choix techniques de BioGate en deux groupes. Les premiers ont été comparés à une autre solution au cours du développement, et la raison du choix est consignée dans le journal de bord et les notes du projet. Les seconds ont été retenus pour leurs avantages, sans comparaison.

\subsection{Choix comparés à une autre solution}

\begin{itemize}
  \item \textbf{FaceNet512 plutôt que VGG-Face} : VGG-Face, associé à une détection de visage basique, s'est révélé trop imprécis sous la lumière d'une webcam. FaceNet512, utilisé avec la bibliothèque DeepFace, a été retenu pour sa robustesse aux conditions lumineuses~\cite{schroff2015facenet,serengil_deepface} ;
  \item \textbf{MediaPipe plutôt que le calcul à la main pour la vivacité} : une première version calculait à la main l'ouverture des yeux et de la bouche. La version actuelle utilise les 52 coefficients d'expression fournis par le modèle de points du visage de MediaPipe~\cite{google_facelandmarker}, ce qui évite ce calcul. Le modèle est publié sous licence Apache~2.0~\cite{mediapipe_blendshape_card} ;
  \item \textbf{Pas d'analyse de la posture} : une analyse posturale en trois dimensions a été étudiée puis abandonnée, car elle était chargée sans jamais être utilisée ;
  \item \textbf{Telegram plutôt que WhatsApp} : l'interface officielle de WhatsApp exige un portefeuille d'entreprise Meta et, depuis le 1\ier{} juillet 2025, facture au message les envois faits par l'entreprise~\cite{meta_wa_overview,meta_wa_pricing}. Un robot Telegram est gratuit par défaut~\cite{telegram_botfaq} et accepte l'envoi d'une photo de 10~Mo au plus~\cite{telegram_botapi} ;
  \item \textbf{BW21-CBV plutôt qu'ESP32-CAM} : la BW21-CBV, envisagée après l'ESP32-CAM, offre un Wi-Fi bi-bande, alors que le module ESP32-CAM, bâti sur l'ESP32, n'émet qu'en 2,4~GHz~\cite{aithinker_bw21_spec,aithinker_esp32cam,espressif_esp32} ;
  \item \textbf{React avec Vite plutôt que Next.js} : Next.js est conçu pour le rendu côté serveur des sites web, ce qui est inutile dans une application de bureau. Il exige un processus Node.js séparé et pose des incompatibilités avec le système de fichiers d'Electron. React avec Vite est plus simple, et son rechargement rapide facilite le développement ;
  \item \textbf{Python 3.11 plutôt que 3.13} : la version 3.13 posait des incompatibilités entre les bibliothèques de calcul et de vision, notamment NumPy et OpenCV. La version 3.11 est compatible avec l'ensemble de celles utilisées.
\end{itemize}

\subsection{Choix retenus pour leurs avantages}

\begin{itemize}
  \item \textbf{FastAPI} : cadre moderne et rapide pour construire des interfaces de programmation en Python, fondé sur les annotations de types standard~\cite{fastapi_docs}. Écrit dans le même langage que les modèles d'intelligence artificielle (YOLO, DeepFace, MediaPipe, OpenCV), il les appelle directement, sans passerelle entre deux langages ;
  \item \textbf{SQLite} : moteur de base de données intégré au programme, sans serveur et sans configuration, qui tient une base complète, avec ses tables et ses index, dans un seul fichier~\cite{sqlite_about}. Il convient à une installation locale unique, se sauvegarde en copiant ce fichier, et son code est dans le domaine public ;
  \item \textbf{YOLOv8 nano} : le plus petit modèle de la famille YOLOv8, avec 3,2 millions de paramètres, entraîné sur la base \gls{coco} qui comprend une classe \og personne \fg{}~\cite{ultralytics_yolov8docs}. Sa légèreté permet de l'exécuter sur un simple processeur, à environ 15 images par seconde ;
  \item \textbf{Suivi par centroïde} : calcul très léger, adapté à un traitement local sur un petit matériel. Il suffit de comparer les centres des boîtes et un histogramme de couleur par individu ;
  \item \textbf{Cascade de Haar} : méthode de détection d'objets efficace, proposée par Viola et Jones et fournie avec OpenCV~\cite{viola2001rapid,opencv_cascade}. Un appel ne coûte que 58 ms sur le matériel d'essai, ce qui permet de filtrer les visages de face avant le calcul lourd ;
  \item \textbf{CLAHE} : égalisation adaptative de l'histogramme qui limite le contraste~\cite{zuiderveld1994clahe,opencv_clahe}. Elle agit par zones de l'image, ce qui rétablit les traits d'un visage en contre-jour ;
  \item \textbf{DeepFace} : bibliothèque publiée sous licence MIT~\cite{serengil_deepface}, qui propose plusieurs modèles de reconnaissance derrière une même interface et gère l'alignement du visage ;
  \item \textbf{Electron} : il permet de maintenir une seule base de code JavaScript, HTML et CSS pour des applications qui fonctionnent sous Windows, macOS et Linux~\cite{electron_docs}. Le tableau de bord tourne ainsi sur n'importe quel ordinateur du réseau local. Electron stocke le jeton chiffré par le système d'exploitation, un installateur se génère en une seule commande, et le flux MJPEG s'affiche avec une simple balise image.
\end{itemize}
Ensemble, ces choix donnent un système local, léger et sans abonnement : chaque composant s'exécute sur un ordinateur ordinaire, et le seul service extérieur du fonctionnement de base est l'envoi de l'alerte.

\addcontentsline{toc}{section}{Conclusion}
\section*{Conclusion}

Ce chapitre a présenté l'analyse des besoins, l'architecture, le pipeline d'analyse, l'authentification multifacteur, le modèle de données, les alertes, les réglages, la gestion des pannes et les choix techniques de BioGate. Il a distingué ce qui est réalisé et testé, comme la caméra, le serveur, le tableau de bord et les alertes, de ce qui est simulé ou prévu, comme le lecteur d'empreinte, les actionneurs et l'écran de la caméra. Le chapitre suivant décrit la mise en œuvre, les essais réalisés et les limites du prototype.
